
\ifintegriereaufgaben
\section*{Anwendungen Aufgaben}\label{MA2-laengen-flaechen-und-volumina-aufgaben}

Erkl�rungen und Beispiele, die zur Bearbeitung der Aufgaben hilfreich sind, finden Sie im Skript\wennreferenzvorhanden{MA2-laengen-flaechen-und-volumina}{ (Abschnitt \abschnittsverweis{MA2-laengen-flaechen-und-volumina})} und, falls Sie nicht weiterkommen, schauen Sie \hyperref[MA2-laengen-flaechen-und-volumina-loesungen]{hier}.

\else
\section{Anwendungen Aufgaben}\label{MA2-laengen-flaechen-und-volumina-aufgaben}

Regeln und Beispiele zu diesen Aufgaben finden Sie im Skript. Die L�sungen finden Sie \hyperref[MA2-laengen-flaechen-und-volumina-loesungen]{hier}.
\fi

\begin{exercisebox}[Kegelvolumen]
	Wie l��t sich ohne Integration beweisen, dass das Volumen eines Kegels mit Grundfl�che \(A\) und H�he \(h\) gerade \(V= \displaystyle{\frac13} A \cdot h\) betr�gt? Hinweis: Benutzen Sie, was Sie �ber das Volumen einer Pyramide mit quadratischer Grundfl�che wissen!
\end{exercisebox}

\begin{exercisebox}[Bogenl�nge] %Papula I, V.10, Aufgabe 15 und 16
	\begin{itemize}
		\item[(a)] Geben Sie das Integral an, das die Bogenl�nge der Kurve \( y = \ln(x^3)\) f�r \( 1 \leq x \leq e\) beschreibt.
		\item[(b)] Geben Sie das Integral an, das die Bogenl�nge der Kurve  \( y = \sqrt{ x^3 } \) f�r \( 1 \leq x \leq 3\) beschreibt.
	\end{itemize}
\end{exercisebox}

\begin{exercisebox}[Fl�cheninhalt zwischen Graphen] %Papula V, Abschnitt 10 Aufgabe 5 und 7
	\begin{itemize}
		\item[(a)] Bestimmen Sie f�r die Funktionen
			\[ f:x\mapsto x^2-2x-1,\quad g:x\mapsto -x+5\]
			den Fl�cheninhalt zwischen den Graphen und ihren Schnittpunkten.
		\item[(b)] Bestimmen Sie die Fl�che \(A\) (!) zwischen der Funktion
			\[ f:x\mapsto \frac15 x ( x^2-4) \]
			und der $x$-Achse im Intervall $x\in [-3,3]$.

	\end{itemize}
\end{exercisebox}


\begin{exercisebox}[Rotationsvolumen] % Papula I, V.10, Aufgabe 13
	Bestimmen Sie das Rotationsvolumen eines K�rpers, der durch Drehung des Kurvenst�cks \(y= \sqrt{ x^2 - 9 }\) f�r \(3\leq x \leq 5\)
	\begin{itemize}
		\item[(a)] um die \(x\)-Achse entsteht.
		\item[(b)] um die \(y\)-Achse entsteht.
	\end{itemize}
\end{exercisebox}

\begin{exercisebox}[Lineare Abbildungen]
Auf einem Schrottplatz werden Autofracks mit einer Presse zusammengestampft. Jede Karosserie wird dabei in jede Richtung gleichm��ig umgeformt und das Volumen, das es urspr�nglich einnahm, wird durch die Presse um $80\%$ verringert.

Geben Sie eine lineare Abbildung an, die den Effekt der Presse simulieren kann (keine eindeutige L�sung).
\end{exercisebox}

%\begin{exercisebox}[Klaus und die Arbeitsplatte]
%	Idee: Arbeitsplatte als 1D Objekt annehmen, ein Ausschnitt f�r das Sp�lbecken. Klaus kann max 30kg tragen, sein Vater 50kg.
%	Wie schwer ist die Arbeitsplatte?
%	Klaus w�rde gerne vorne tragen, kann das klappen?
%\ende{exercisebox}

%\begin{exercisebox}[Kugelvolumen]
%	Zeigen Sie ohne Integration, dass das Volumen einer Kugel mit Radius \(r\)
%		\[
%			V = \frac43\pi r^3
%		\]
%		betr�gt.
%
%%\vspace{0.3cm}
%%{\bf L�sung:} \newline
%	%Eine Kugel mit Radius $r$ hat Schicht f�r Schicht die gleiche Querschnittsfl�che wie ein gerader Kreiszylinder mit Radius $r$ und H�he $2r$, aus dem man oben und unten einen Kegel herausschneidet.
%
%		%Wenn man zeigen kann, dass zwei K�rper dieselbe H�he und in jeder H�he dieselbe Querschnittfl�che hat, dann haben die zwei K�rper auch dasselbe Volumen.
%
%	\end{exercisebox}

%\subsection{Praxis}

%\begin{exercisebox}[Kurvenintegral]
%	Berechnen Sie die Kurvenintegrale der gegebenen Sklarfelder $\rho:\Gamma \to \mathbb R$ l�ngs der Kurve $\Gamma$:
%	\begin{itemize}
%		\item[(a)] $\rho(x,y) = x+y, \quad \Gamma: t\mapsto (\cos(t), \sin(t))^\intercal, \; t\in [0,\pi]\,.$
%		\item[(b)] $\rho(x,y) = \displaystyle{\frac{x}{x^2+y^2}}, \quad \Gamma: t\mapsto (\cos(t), \sin(t))^\intercal, \; t\in [-\displaystyle{\frac{\pi}{2}},\displaystyle{\frac{\pi}{2}}]\,.$
%		\end{itemize}
%\vspace{0.3cm}
%{\bf L�sung:} \newline
%\begin{itemize}
%	\item[(a)]
%		\[ \int\limits_\Gamma \rho(x,y)\,\mathrm{d}\boldsymbol s = \int\limits_\Gamma \rho(\Gamma(t)) \, \mathrm{d}\boldsymbol s(t) = \int\limits_0^\pi \rho(\Gamma(t))\cdot |\dot \Gamma(t)| \, \mathrm{d}t\]
%		Die Parametrisierung des Bogens durch $\Gamma$ ist stetig differenzierbar, also
%		\[ \mathrm{d}\boldsymbol s(t) = |\dot \Gamma(t)| \mathrm{d}t = \sqrt{\sin^2(t) + \cos^2(t)}\,\mathrm{d}t = \mathrm{d}t\,. \]
%		Und damit lautet das Kurvenintegral:
%		\[ \int\limits_\Gamm \rho(x,y)\,\mathrm{d}\bolsymbol s = \int\limits_0^\pi (\cos(t)+ \sin(t)) \, \mathrm{d}t = \Big[ \sin(t) - \cos(t)\Big]_0^\pi = 1+1 = 2\,.
%		\]
%	\item[(b)]
%		\[ \int\limits_\Gamma \rho(x,y)\,\mathrm{d}\boldsymbol s = \int\limits_\Gamma \frac{x}{x^2+y^2} \, \mathrm{d}\boldsymbol s\]
%		Die Parametrisierung des Bogens durch $\Gamma$ ist stetig differenzierbar, also  $x=\cos(t)$, $y=\sin(t)$ und $\mathrm{d}\boldsymbol s(t)  \mathrm{d}t$. Und damit lautet das Kurvenintegral:
%		\[ \int\limits_\Gamm \rho(x,y)\,\mathrm{d}\bolsymbol s = \int\limits_{-\pi/2}^{\pi/2}  \frac{\cos(t)}{ \cos^2(t) + \sin^2(t)} \, \mathrm{d}t = \\int\limits_{-\pi/2}^{\pi/2}  \cos(t) \, \mathrm{d}t = \Big[ \sin(t) \Big]_{-\pi/2}^{\pi/2} = 2\,.
%		\]
	%\end{exercisebox}

%Moodle %HMI I, Blatt 11, Aufgabe 1 Moodle Polygonzug

